\setcounter{section}{1}
\setcounter{thm}{0}
Let us formally state our problems and results.
%
\begin{problem}
Let $\confam$ be a finite \emph{constraint family} of Boolean functions. The problem $\SAT(\confam)$ asks to determine whether a given formula $\phi$ on Boolean variables $x_1,\dots,x_n$ is satisfiable by an assignment with $k$ ones, where $\phi$ is a conjunction of $m$ constraints $C$ of the form $f(\mathbf{x})$, where $f:\{0,1\}^r \to \{0,1\}$ is a constraint function in $\confam$ and $\mathbf{x}$ is an $r$-tuple of variables among $x_1,\dots, x_n$.
\end{problem}
%
Note that if all $f\in \confam$ have arity bounded by $r$, then there are at most $\Oh(n^r)$ possible constraints, and exhaustive search solves $\SAT(\confam)$ in time $\Oh(n^{k+r})$.
%
We will show that the complexity of $\SAT(\confam)$ is tightly characterized by the set of functions \emph{expressible as restrictions of constraint functions $f\in \confam$}. To formally introduce this concept, let $f: \{0,1\}^r \to \{0,1\}$ be an arbitrary Boolean function. We say that $g:\{0,1\}^s\to \{0,1\}$ is a \emph{restriction of $f$} if it is obtained from $g$ by replacing each argument of $f$ by either the constant $0$, the constant $1$, or an argument of $g$, i.e., we can partition $[r]$ into $X_1,\dots, X_s, Z_0, Z_1$ such that 
\[ g(x_1,\dots,x_s) = f(\overbrace{\entry{x_1}}^{X_1}, \dots, \overbrace{\entry{x_s}}^{X_s}, \overbrace{\entry{0}}^{Z_0}, \overbrace{\entry{1}}^{Z_1}).\]
Here, $\overbrace{\entry{y}}^Y$ denotes plugging in $y$ for all (not necessarily contiguous) positions $Y\subseteq [r]$, see Section~\ref{sec:prelim}. 
%
   
\begin{defn}
Let $g:\{0,1\}^d\to \{0,1\}$ be an arbitrary Boolean function.
A constraint family $\confam$ \emph{represents} $g$ if there is some $f\in \confam$ such that $g$ is a restriction of $f$. If $\confam$ does not represent~$g$, we say that $\confam$ \emph{avoids} $g$.
\end{defn}
%
Let $\Impl:\{0,1\}^2 \to \{0,1\}$ and $\NAND_d :\{0,1\}^d \to \{0,1\}$  be the binary implication and $d$-ary $\NAND$ function, respectively, i.e., 
\begin{align*}
\Impl(y_1,y_2) & \coloneqq \overline{y_1} \vee y_2,\\
\NAND_d(y_1,\dots,y_d) & \coloneqq \overline{\bigwedge\nolimits_{i=1}^d y_i}.  
\end{align*}
%
In~\cite{Marx05}, it is shown that assuming $\FPT\ne\W[1]$, $\SAT(\confam)$ is solvable in FPT time $f(k)n^{c}$ if and only if $\confam$ is \emph{weakly separable}, which is a condition equivalent to $\confam$ avoiding $\NAND_2$ and $\Impl$.
We show an almost tight characterization of $g(k)$ (under plausible assumptions from fine-grained complexity theory) that depends only on whether or not $\confam$ represents $\Impl$, $\NAND_2$ or $\NAND_d$ for higher order $d\ge 3$. Specifically, we obtain the following main theorem, illustrated in Figure~\ref{fig:results}.
%
\begin{figure}
\centering
\includegraphics[width=0.85\textwidth]{ccc2020-27-classification-overview.pdf}
\caption{Overview over our main results. The parts of the diagram to the right of the vertical $\Impl$ line depict $\confam$ representing $\Impl$, while the parts to the left avoid $\Impl$. Analogously, the parts of the diagram above a $\NAND_d$ line depict $\NAND_d$-representing $\confam$, while those below avoid $\NAND_d$. For each cell, we illustrate our (typically matching) algorithmic and hardness results, together with a problem that is complete for this cell (in a certain sense). For clarity of presentation, we drop additional $f(k)n^c$-factors of stated running times.}
\label{fig:results}
\end{figure}
%
%
\newcommand{\regime}[1]{{\textbf{\textup{[#1]}}}}
\begin{thm}
Let $\confam$ be a finite constraint family.
\begin{enumerate}
\item \regime{FPT regime} If $\confam$ avoids both $\NAND_2$ and $\Impl$, then there is a computable $f(k)$ and constant~$c_\confam$ such that $\SAT(\confam)$ can be solved in time $f(k)n^{c_\confam}$.
\item \regime{Subexponential regime}\newline If $\confam$ represents $\Impl$, but avoids $\NAND_2$, then there is a computable $f(k)$ and constant~$c_\confam$ such that $\SAT(\confam)$ can be solved in time $f(k)n^{4\sqrt{k}+c_\confam}$; \newline furthermore, for no computable $f(k)$ and constants $c_\confam, \eps > 0$, $\SAT(\confam)$ can be solved in time  $f(k)n^{(\omega/6-\eps)\sqrt[3]{k}+c_\confam}$, unless the $k$-clique conjecture fails.
\item \regime{Clique regime}\newline  If $\confam$ represents $\NAND_2$, but avoids $\NAND_3$, then there is a computable $f(k)$ and constant~$c_\confam$ such that $\SAT(\confam)$ can be solved in time $f(k) n^{(\omega/3)k+c_\confam}$; \newline furthermore for no computable $f(k)$ and constants $c_\confam, \eps > 0$, $\SAT(\confam)$ can be solved in time $f(k) n^{(\omega/3-\eps)k+c_\confam}$, unless the $k$-clique conjecture fails. 
\item \regime{Brute-force regime}\newline  If $\confam$ represents $\NAND_3$, then for no computable $f(k)$ and constants $c_\confam, \eps > 0$, $\SAT(\confam)$ can be solved in time $f(k)n^{(1-\eps)k+c_\confam}$, unless the 3-uniform $k$-HyperClique conjecture fails. 
\end{enumerate}
\end{thm}
%
%
That is, we only have four regimes: $g(k)$ is either constant, sublinear in $k$ with a value between essentially $(\omega/6) \sqrt[3]{k}$ and $4\sqrt{k}$, the clique detection bound of essentially $(\omega/3) k$, or the brute force bound of essentially $k$.  Note that we do not try to optimize the bounds on $f(k)$, which generally are bounded by $r^{\Oh(k^3)}$, where $r$ is the arity of $\confam$.
%
